\documentclass[11pt]{article}
\usepackage[margin=1in, headheight=14pt]{geometry}
\usepackage{amsmath, amssymb}
\usepackage{hyperref}
\usepackage{enumitem}
\usepackage{algorithm}
\usepackage{algpseudocode}
\usepackage{fancyhdr}
\usepackage{tcolorbox}
\usepackage{microtype}
\usepackage{tikz}
\usepackage{pgfplots}
\pgfplotsset{compat=1.18}
\usetikzlibrary{arrows.meta, positioning, calc, shapes.geometric}

\input{notation.tex}
\input{zk-notation.tex}
\input{environments.tex}
\input{lecture-macros.tex}

\pagestyle{fancy}
\fancyhf{}
\lhead{EECS 498/CSE 598: Zero-Knowledge Proofs}
\rhead{Winter 2026}
\lfoot{Lecture 15: ASIC-Model Frontends II: Programs in R1CS}
\rfoot{\thepage}
\renewcommand{\headrulewidth}{0.4pt}
\renewcommand{\footrulewidth}{0.4pt}

% Lecture-specific macros
\tcbuselibrary{skins, breakable}

\newtcolorbox{graybox}[1][]{
  colback=gray!8,
  colframe=gray!50,
  fonttitle=\bfseries,
  #1
}

\begin{document}

%-----------------------------------------------------------------------
% Title Block
%-----------------------------------------------------------------------
\begin{center}
  {\LARGE \textbf{EECS 498 / CSE 598: Zero-Knowledge Proofs}}\\[6pt]
  {\large Winter 2026}\\[4pt]
  {\large \textbf{Lecture 15: ASIC-Model Frontends II: Programs in R1CS}}\\[4pt]
  {\large Instructor: Paul Grubbs \quad
    \href{mailto:paulgrub@umich.edu}{\texttt{paulgrub@umich.edu}} \quad
    Beyster 4709}
\end{center}

\vspace{0.5em}
\hrule
\vspace{1em}

%-----------------------------------------------------------------------
% Agenda
%-----------------------------------------------------------------------
\section*{Agenda}
\begin{itemize}
  \item ASIC-model frontends and the $\algo{R1CS}$ target
  \item Encoding basic arithmetic programs
  \item Conditionals, bit checks, and range checks
  \item Boolean and integer operations
  \item Optimization themes
\end{itemize}

\hrule
\vspace{1em}

%-----------------------------------------------------------------------
\section{ASIC-Model Frontends and \texorpdfstring{$\algo{R1CS}$}{R1CS}}
%-----------------------------------------------------------------------

ASIC-model frontends target an algebraic constraint language directly.  One
common target is a rank-$1$ constraint system.

\begin{definition}[Rank-$1$ constraint system]
Fix a witness vector $z$.  A rank-$1$ constraint is an equation of the form
\[
  \langle a, z \rangle \cdot \langle b, z \rangle = \langle c, z \rangle,
\]
where $a$, $b$, and $c$ are coefficient vectors over the backend field.  An
\algo{R1CS} instance is a collection of such constraints, equivalently written
as
\[
  Az \circ Bz = Cz.
\]
\end{definition}

The vector $z$ usually contains a distinguished constant wire equal to $1$
together with the ordinary program variables and any auxiliary variables
introduced during compilation.

The source languages feeding these systems often look closer to hardware or
constraint descriptions than to full general-purpose software.  Hardware-style
languages such as \algo{Verilog} and \algo{VHDL} provide a useful point of
comparison because the compilation target is again a fixed dependency graph.

%-----------------------------------------------------------------------
\section{Basic Program Constructs}
%-----------------------------------------------------------------------

\subsection{Addition and Multiplication}

The simplest program fragments map directly to one constraint each.  For
example,
\[
\begin{array}{l}
\algo{add1}(x): \\
\quad y \leftarrow x + 1
\end{array}
\]
becomes
\[
  1 + x - y = 0.
\]

Likewise,
\[
\begin{array}{l}
\algo{add}(a,b): \\
\quad c \leftarrow a+b
\end{array}
\qquad\leadsto\qquad
a+b-c=0
\]
and
\[
\begin{array}{l}
\algo{mul}(a,b): \\
\quad c \leftarrow ab
\end{array}
\qquad\leadsto\qquad
ab=c.
\]

The general compilation pattern is to introduce a new variable for every
intermediate value and then constrain that variable by a low-degree equation.

\subsection{Composite Arithmetic Expressions}

Consider
\[
\begin{array}{l}
\algo{addmul}(a_0,a_1,b_0,b_1): \\
\quad \text{return } a_0b_0 + a_1b_1.
\end{array}
\]
A standard static-single-assignment decomposition is
\[
  c_0 = a_0b_0,\qquad
  c_1 = a_1b_1,\qquad
  d = c_0 + c_1.
\]
The resulting \algo{R1CS} description uses one multiplication constraint for
each $c_i$ and one linear constraint for the final addition.

Division is represented by rewriting it as a multiplicative relation.  For
\[
\begin{array}{l}
\algo{divide}(a,b): \\
\quad c \leftarrow a/b,
\end{array}
\]
the checked form is
\[
  cb = a.
\]
If the intended source semantics excludes division by zero, then the frontend
must also enforce the side condition $b \neq 0$ by an auxiliary gadget.

\begin{proposition}
Straightline arithmetic over a field can be compiled into \algo{R1CS} by
introducing variables for intermediate values and constraining each arithmetic
step by addition and multiplication relations.
\end{proposition}

\begin{proof}[Proof sketch]
Every straightline program is an acyclic expression graph.  Replace each gate
output by a fresh variable.  Addition gates contribute linear constraints and
multiplication gates contribute rank-$1$ constraints.  The final output
variable is then constrained to equal the source-level result.
\end{proof}

%-----------------------------------------------------------------------
\section{Conditionals, Bits, and Ranges}
%-----------------------------------------------------------------------

Conditionals are not native field operations, so the control bit must itself be
represented inside the constraint system.

\begin{graybox}[title={Conditional-Multiplication Gadget}]
For bit $b \in \{0,1\}$ and branch values $xy$ and $xz$, the conditional
product
\[
  \algo{condMul}(x,y,z,b)
  =
  \begin{cases}
    xy & \text{if } b=1, \\
    xz & \text{if } b=0
  \end{cases}
\]
can be encoded with auxiliary variables
\[
  c_1 = xy,\qquad
  b' = 1-b,\qquad
  c_0 = xz,
\]
\[
  d_0 = bc_1,\qquad
  d_1 = b'c_0,\qquad
  d = d_0 + d_1.
\]
The output is $d$ together with the bitness constraint
\[
  b(b-1)=0.
\]
\end{graybox}

Bit checks and range checks are the other indispensable gadgets.

\begin{definition}[Bit check]
A field element $b$ is constrained to be Boolean by the equation
\[
  b(b-1)=0.
\]
\end{definition}

\begin{proposition}[Range check by bit decomposition]
To prove that
\[
  x \in [0,7],
\]
it is enough to introduce bits $b_0,b_1,b_2$ and enforce
\[
  x = 4b_2 + 2b_1 + b_0
\]
together with the three bitness constraints
\[
  b_i(b_i-1)=0
  \qquad\text{for } i \in \{0,1,2\}.
\]
\end{proposition}

\begin{proof}
Every value in $[0,7]$ has a unique binary expansion using three bits, and the
three bitness constraints ensure that no other values are possible.
\end{proof}

%-----------------------------------------------------------------------
\section{Boolean and Integer Operations}
%-----------------------------------------------------------------------

Once variables are known to be bits, standard Boolean operations become
low-degree polynomials:
\[
  a \land b = ab,
  \qquad
  a \lor b = a + b - ab,
  \qquad
  \lnot a = 1-a,
\]
\[
  a \oplus b = a + b - 2ab.
\]

These formulas are the reason bit decomposition is so useful.  Integer
operations can often be reduced to bit-level constraints plus carry or borrow
relations.  For example, a less-than test
\[
  \algo{lessthan}(x,y)
\]
is usually compiled by decomposing $x$ and $y$ into bits and checking the first
position at which they differ.

\begin{remark}
The native field arithmetic is efficient; the expense comes from recovering the
intended integer or bit-string semantics inside that field.
\end{remark}

%-----------------------------------------------------------------------
\section{Optimization Themes}
%-----------------------------------------------------------------------

Constraint efficiency depends less on the source syntax than on the choice of
representation.  The main recurring levers are:
\begin{itemize}
  \item reuse intermediate variables instead of recomputing them,
  \item minimize the number of multiplication constraints,
  \item reduce the number of range checks and bit decompositions,
  \item exploit backend-specific features such as lookups or custom gates when
        the target model supports them.
\end{itemize}

The common pattern is the same one already seen in frontend design more
generally: a mathematically equivalent encoding may still be far more expensive
for the proof backend.

\end{document}
